\chapter{Project Management and Process}
\label{ch:project-management-and-process}
\section{Timeline}
The project ran over the course of the six-month internship, organized into two main phases rather than a fixed day-by-day schedule. The initial build phase established the system's foundation, in order: source selection and scope, schema design, parsing, normalization and validation, node construction, knowledge graph construction, the NLP/search layer, the agent, and the web backend. This was followed by a substantially longer second phase — driven by direct feedback from both the intern and the internship supervisors after the initial system was demonstrated — that added the hybrid LLM layer, grounded generation, the analytics dashboard, the React frontend rewrite, containerized deployment, hosting, multi-language support, and the four chatbot-quality features detailed in Section 10.6. In total, the system took a bit more than four months of active development to reach the state described in this report, with the iterative, decision-log-driven working method described below applied consistently across both phases.

\section{Decision-logging discipline}
Every scope decision, every confirmed bug, and every accepted trade-off in this project is recorded, dated, in a running decision log (\texttt{docs/\allowbreak{}decisions.md}) specifically so the reasoning behind the current system is traceable without needing to reconstruct it from memory or chat history — the primary source for this report itself, and the artifact this report's author considers the single most valuable piece of documentation the project produced, independent of the code.

\section{Working method}
Two working principles were applied consistently enough, across a genuinely large number of individually-documented bugs, to be considered defining characteristics of how this project was executed rather than incidental good habits: first, no fix was trusted until it had been checked against real data — a real screenshot, real page characters, a real confirmed user question, or a whole-document consistency scan — rather than reasoned about in the abstract; second, when a fix attempt was tried and found not to work, or found to introduce a new regression, it was reverted and documented as an honest, explained dead end rather than kept because it looked plausible or silently discarded without a record. Section 9.4's account of the abandoned parsing fix during the grounding-fallback investigation is one concrete instance of the second principle in practice, not the only one.

